\section*{Introdução}
O trabalho visa resolver um problema de classificação binária levantado pelo Jardim Comunitário Raízes Vivas, focando no uso de técnicas de Machine Learning e pré-processamento de dados.
O objetivo é desenvolver um modelo preditivo capaz de determinar o \textbf{`support\_outcome`} para membros do programa, classificando-o como "bom" (good) ou "ruim" (bad).

Para realizar esta tarefa, utilizamos um conjunto de dados de treinamento fornecido no Moodle (800 amostras e 21 atributos), detalhando o perfil demográfico, financeiro e histórico de participação de cada membro.
Este relatório detalha a metodologia de pré-processamento, a seleção de atributos e a subsequente implementação e ajuste de seis modelos de classificação distintos:
\begin{itemize}
    \item Árvore de Decisão (Decision Tree)
    \item Floresta Aleatória (Random Forest)
    \item Aumento de Gradiente (Gradient Boosting)
    \item Classificador Bayesiano (Gaussian Naive Bayes)
    \item Máquina de Vetor de Suporte (Support Vector Classifier)
    \item Classificador de Votação (Voting Classifier)
\end{itemize}
Os modelos são avaliados usando validação cruzada (K-Fold Cross-Validation com k=5) e comparados com base nas métrica acurácia e a AUC de cada modelo. Para encontrar o melhor modelo de cada algoritmo, a técnica de grid-search foi aplicada para realizar o tuning de hiperparâmetros.